\section{Independent audit of successor-only uniqueness}
The proposed argument from successor-only uniqueness (SU) to contextual
zero-equation induction is correct in ordinary intuitionistic first-order
logic with actual equality, scalar computation, numerical coverage, and
$0\ne S0$. One minor derived step should be made explicit: with
$Z(v)=I(v,0;u.S0)$, an equation $Sx=0$ implies
$S0=Z(Sx)=Z(0)=0$. Thus every successor is apart from zero.
Coverage then makes each zero equation $q(n)=0$ decidable, justifying the
finite case split in the retraction construction without any general
excluded-middle principle.

The retraction proof uses no induction before SU. In its successor case,
$e(r(Sm))=r(Sm)$ follows from idempotence of $e$ at
$S(e(r(m)))$. The equation $r(Sn)=S(r(n))$ then follows from
$e(S(e(x)))=S(e(x))$ and the already established fixed-point equation.
Applying SU to $r$ and separately to the identity function proves
$r(n)=n$. All arguments allow a fixed retained context $\Gamma$.

\paragraph{The existing zero-induction countermodel does not refute SU.}
In the two-point-fiber model of the repository's
\texttt{canonical\_recursors} zero-induction counterexample, define
\[
 q(n)=I(n,n;u.u).
\]
For standard $n$ this equals $n$. For nonstandard $n=(m,i)$, the specified
identity-iterator choice gives $q(n)=(m,0)$. Consequently
\[
 q(0)=0,\qquad q(Sn)=S(q(n))
\]
hold everywhere. SU for $q$ would give
$q(n)=I(n,0;u.Su)$, whereas SU for $f(n)=n$ would give
$n=I(n,0;u.Su)$. Their conjunction contradicts a nonstandard label-one
point. This conclusion is independent of the model's choices of the
successor iterator: those two SU instances cannot both hold.
The earlier fiber countermodel therefore separates zero induction from
full equation induction, but it cannot settle whether SU is already
sufficient. No model of SU refuting full equation induction is established
by this audit.

\paragraph{Orbit needs counter independence.}
When the update is read from a proposed clause or closure, its only active
variable must be the previous state. It must not additionally refer to the
outer counter. Without that check, $f(y)=y$ and $h(y,u)=Sy$ satisfy the
proposed local recurrence but give the false orbit conclusion
$1=I(1,0;u.S1)=2$. The delivered checker explicitly rejects that coordinate,
including occurrences under nested iterator binders.
